% Fichier généré automatiquement par tools/generate_corrections.py.
% Source : DYN/DYN-05-Methode/46_RR_RSG/46_RR_RSG.tex
% Statut : complété (3/3 question(s)).
\begin{corrigebox}[Corrigé complété]
\CorrigeQuestion{1}
\begin{marginfigure}
\includegraphics[width=\linewidth]{46_01_cor}
\end{marginfigure}
\CorrigeQuestion{2}
\begin{itemize}
\item Première équation : 
\begin{itemize}
\item On isole 2.
\item Bilan des actions mécaniques extérieures :
\begin{itemize}
\item liaison pivot en $A$ telle que $\vectm{A}{1}{2}\cdot \vect{k_0} = \vect{0}$;
\item pesanteur en $B$ : $\torseurstat{T}{\text{pes}}{2} = \torseurl{-m_2 g \vect{j_0}}{\vect{0}}{B}$;
\item couple moteur : $\torseurstat{T}{1_m}{2} = \torseurl{\vect{0}}{C_m\vect{k_0}}{\forall P}$.
\end{itemize}
\item On applique le théorème du moment dynamique en $A$ en projection sur $\vect{k_0}$ :
$\vectmd{A}{2}{0}\cdot \vk{0} = 0 + \left(\vect{AG_2}\wedge -m_2 g \vect{j_0}\right)\cdot \vk{0} + C_m$. 

\end{itemize}
\item Deuxième équation : 
\begin{itemize}
\item On isole 1+2.
\item Bilan des actions mécaniques extérieures :
\begin{itemize}
\item liaison ponctuelle avec RSG en $I$ telle que $\vectm{I}{0}{1}\cdot \vect{k_0} = \vect{0}$; 
\item pesanteur en $G_1$  : $\torseurstat{T}{\text{pes}}{1} = \torseurl{-m_1 g \vect{j_0}}{\vect{0}}{G_1}$;
\item pesanteur en $G_2$  : $\torseurstat{T}{\text{pes}}{2} = \torseurl{-m_2 g \vect{j_0}}{\vect{0}}{G_2}$.
\end{itemize}
\item On applique le théorème du moment dynamique en $I$ en projection sur $\vect{k_0}$ :
$\vectmd{I}{1+2}{0}\cdot \vk{0} = 0 + \left(\vect{IG_2}\wedge -m_2 g \vect{j_0}\right)\cdot \vk{0} +\left(\vect{IG_1}\wedge -m_1 g \vect{j_0}\right)\cdot \vk{0} $. 
\end{itemize}
\end{itemize}
\marginnote{Remarque : on ne modélise pas la résistance au roulement.}
\CorrigeQuestion{3}
On choisit les inconnues et les conventions positives de la figure,
        on écrit les relations de conservation/fermeture adaptées, puis on résout le
        système obtenu. Le résultat symbolique doit être homogène et vérifié dans les cas
        limites (entrée nulle, rapport unitaire ou position de référence).
\end{corrigebox}
